% -----------------------------------------------------------
\section{Life, Live}
% -----------------------------------------------------------
	\label{LIFE}
	
% % ----------------------
% \subsection{See also}
% % ----------------------

% % ----------------------
% \subsection{1828 American Dictionary of the English Language} % *1828
% % ----------------------
	% \label{LIFE-1828}

	% \begin{compactitem}
		% \item
	% \end{compactitem}

% % ----------------------
% \subsection{Oxford English Dictionary} % *OED
% % ----------------------
	% \label{LIFE-OED}

	% \begin{compactitem}
		% \item[]
	% \end{compactitem}

% ----------------------
\subsection{Scriptural Usage} % *SCRIPTURES
% ----------------------
	\label{LIFE-SCRIPTURES}

	% % -----
	% \subsubsection{Old Testament} % *OT
	% % -----
		% \label{LIFE-OT}
	
		% % --
		% \paragraph{KJV}
		% % --
			% \label{LIFE-OT-KJV}
	
			% \begin{tabular}{ll}
				% Total occurrences in the OT & 
			% \end{tabular}
			
			% \begin{compactdesc}
				% \item[]
			% \end{compactdesc}
			
		% % --
		% \paragraph{Hebrew Bible}
		% % --
			% \label{LIFE-HEBREW}
		
			% %
			% \subparagraph{\texorpdfstring{\he{sample word}}{}}
			% %
				% \label{PAGA} % whatever is in step bible without periods
			
				% \begin{compactdesc}
					% \item[HALOT , PDF ] \textbf{}
						% \begin{compactitem}
							% \item[1]
						% \end{compactitem}
					% \item[BDB , PDF ] \textbf{}
						% \begin{compactitem}
							% \item[1]
						% \end{compactitem}
					% \item[DCH PDF ] \textbf{}
						% \begin{compactitem}
							% \item[1]
						% \end{compactitem}
				% \end{compactdesc}
				
				% \begin{compactdesc}
					% \item[Some OT uses] \textbf{}
						% \begin{compactdesc}
							% \item[]
						% \end{compactdesc}
				% \end{compactdesc}
			
		% % --
		% \paragraph{Septuagint}
		% % --
			% \label{LIFE-LXX}
		
			% %
			% \subparagraph{\texorpdfstring{\gr{sample word}}{}}
			% %
		
	% -----
	\subsubsection{New Testament} % *NT
	% -----
		\label{LIFE-NT}
	
		% --
		\paragraph{KJV}
		% --
			\label{LIFE-NT-KJV}
			
	
			% \begin{tabular}{ll}
				% Total occurrences in the NT & 
			% \end{tabular}
			
			\begin{compactdesc}
				\item[{\hyperref[JOHN-5-26]{John 5:26}}]
			\end{compactdesc}
			
		% --
		\paragraph{Greek New Testament} % *GREEK
		% --
			\label{LIFE-GREEK}
		
			%
			\subparagraph{\texorpdfstring{\gr{ζάω}}{}}
			%
				\label{ZAO}
			
				\begin{compactdesc}
					\item[BDAG 375a] \textbf{}
						\begin{compactitem}
							\item[1] to be alive physically, \textit{live}
							\item[1a] of physical life in contrast to death
							\item[1a$\alpha$] generally
							\item[1a$\beta$] of dead persons who return to life \textit{become alive again}
							\item[1a$\gamma$] of sick persons, if their illness terminates not in death but in recovery \textit{be well, recover}
							\item[1a$\delta$] also of healthy persons, \textit{live on, remain alive}
							\item[1a$\epsilon$] of beings that in reality, or as they are portrayed, are not subject to death
							\item[1b] with mention of that upon which life depends
							\item[1c] with more precise mention of the sphere
							\item[2] to live in a transcendent sense, \textit{live}, of the sanctified life of a child of God
							\item[2a] in the world
							\item[2b] in the glory of the life to come
							\item[2b$\beta$] more specifically \gr{εἰς τὸν αἰῶνα} \textit{have eternal life}
							\item[3] to conduct oneself in a pattern of behavior, \textit{live}
							\item[3a] used with adverbs or other modifiers
							\item[3b] \textit{live for someone or something, for the other's benefit}
							\item[4] to be full of vitality, \textit{be lively}
							\item[5] to be life-productive, \textit{offer life} participle used with respect to things
						\end{compactitem}
					\item[CGL , PDF ] \textbf{}
						\begin{compactitem}
							\item 
						\end{compactitem}
					\item[LSJ , PDF ] \textbf{}
						\begin{compactitem}
							\item 
						\end{compactitem}
				\end{compactdesc}
				
				\begin{compactdesc}
					\item[Some NT uses] \textbf{}
						\begin{compactdesc}
							\item[{\hyperref[JOHN-4-10]{John 4:10}}]
						\end{compactdesc}
				\end{compactdesc}
				
			%
			\subparagraph{\texorpdfstring{\gr{ψυχή}}{}}
			%
				\label{PSUCHE}
				
				% \begin{tabular}{ll}
					% Total occurrences in the NT & 
				% \end{tabular}
				
				\begin{compactdesc}
					\item[Some NT uses] \textbf{} % Change to "all" if they're all listed
						\begin{compactdesc}
							\item[{\hyperref[MATTHEW-10-39]{Matthew 10:39}}]
						\end{compactdesc}
				\end{compactdesc}
			
				\begin{compactdesc}
				
					\item[BDAG 979a] It is oft. impossible to draw hard and fast lines in the use of this multivalent word. Generally it is used in reference to dematerialized existence or being, but, apart from other data, the fact that \gr{ψυχή} is also a dog's name suggests that the primary component is not metaphysical...Without \gr{ψυχή} a being, whether human or animal, consists merely of flesh and bones and without functioning capability. Speculations and views respecting the fortunes of \gr{ψυχή} and its relation to the body find varied expression in our literature.
						\begin{compactitem}
							\item[1] \textbf{life on earth in its animating aspect making bodily function possible}
							\item[1.a] \textbf{\textit{(breath of) life, life-principle, soul}}...when it leaves the body death occurs \hyperref[LUKE-12-20]{Luke 12:20}...The soul is delivered up to death (the pass. in reference to divine initiative), i.e. into a condition in which it no longer makes contact with the physical structure it inhabited 1 Cl 16:13 (\hyperref[ISAIAH-53-12]{Is 53:12}), whereupon it leaves the realm of earth and lives on in Hades...\hyperref[ACTS-2-27]{Ac. 2:27} (\hyperref[PSALMS-15-10]{Ps 15:10}), \hyperref[ACTS-2-31]{31} v.l. or some other place outside the earth \hyperref[REVELATIONNT-6-9]{Revelation 6:9}; \hyperref[REVELATIONNT-20-4]{20:4}
							\item[1.b] the condition of being alive, \textbf{\textit{earthly life, life}} itself
							\item[1.c] by metonymy, \textbf{that which possesses life/soul}
							\item[2] \textbf{seat and center of the inner human life in its many and varied aspects, \textit{soul}}
							\item[2.a] of the desire for luxurious living...the soul as the seat of enjoyment of the good things in life...
							\item[2.b] of evil desires
							\item[2.c] of feelings and emotions
							\item[2.d] as the seat and center of life that transcends the earthly...Opp. Tat. 13, 1, who argues the state of the \gr{ψυχή} before the final judgment and states that it is not immortal per se but experiences the fate of the body \gr{οὐκ ἔστιν ἀθάνατος}. As such it can receive divine salvation \gr{σῴζου σὺ καὶ ἡ ψυχή σου} \textit{be saved, you and your soul}...Humans cannot injure it, but God can hand it over to destruction \hyperref[MATTHEW-10-28]{Matthew 10:28ab}...There is nothing more precious than \gr{ψυχή} in this sense \hyperref[MATTHEW-16-26]{Mt 16:26b}; \hyperref[MARK-8-37]{Mk 8:37}. It stands in contrast to \gr{σῶμα}, in so far as that is \gr{σάρξ}...
							\item[2.e] Since the soul is the center of both the earthly (1a) and the transcendent (2d) life, pers. can find themselves facing the question concerning the wish to ensure it for themselves...The contrast betw. \gr{τὴν ψυχὴν εὑρεῖν} and \gr{ἀπολέσαι} is found in \hyperref[MATTHEW-10-39]{Mt 10:39ab}...
							\item[2.f] On the combination of \gr{ψυχή} and \gr{πνεῦμα} in \hyperref[1THESSALONIANS-5-23]{1Th 5:23}; \hyperref[HEBREWS-4-12]{Hb 4:12}...
							\item[2.g] In var. Semitic languages the reflexive relationship is paraphrased with \he{nEpE+s}...the corresp. use of \gr{ψυχή} may be detected in certain passages in our lit., esp. in quots. fr. the OT and in places where OT modes of expr. have had considerable influence...
							\item[3] \textbf{an entity with personhood, \textit{person}}
						\end{compactitem}
						
					% \item[Brill , PDF ] \textbf{}
						% \begin{compactitem}
							% \item 
						% \end{compactitem}
						
					% \item[CGL , PDF ] \textbf{}
						% \begin{compactitem}
							% \item 
						% \end{compactitem}
						
					% \item[LSJ , PDF ] \textbf{}
						% \begin{compactitem}
							% \item 
						% \end{compactitem}
						
					% \item[TDNT ] \textbf{}
						% \begin{compactitem}
							% \item 
						% \end{compactitem}
						
					\item[TDNTA PDF 1221--1230] \textbf{}
						\begin{compactitem}
							\item 
						\end{compactitem}
						
					% \item[NIDNTT ] \textbf{}
						% \begin{compactitem}
							% \item 
						% \end{compactitem}
						
					% \item[NIDNTTE ] \textbf{}
						% \begin{compactitem}
							% \item 
						% \end{compactitem}
						
				\end{compactdesc}
		
	% % -----
	% \subsubsection{Book of Mormon} % *BOM
	% % -----
		% \label{LIFE-BOM}
	
		% \begin{tabular}{ll}
			% Total occurrences in the BOM & 
		% \end{tabular}
		
		% \begin{compactdesc}
			% \item[]
		% \end{compactdesc}
		
	% -----
	\subsubsection{Doctrine and Covenants} % *DC
	% -----
		\label{LIFE-DC}
	
		% \begin{tabular}{ll}
			% Total occurrences in the DC & 
		% \end{tabular}
		
		\begin{compactdesc}
			\item[{\hyperref[DC66-2]{DC 66:2}}]
		\end{compactdesc}
		
	% % -----
	% \subsubsection{Pearl of Great Price} % *PGP
	% % -----
		% \label{LIFE-PGP}
	
		% \begin{tabular}{ll}
			% Total occurrences in the PGP & 
		% \end{tabular}
		
		% \begin{compactdesc}
			% \item[]
		% \end{compactdesc}
		
% % ----------------------
% \subsection{Key Phrases} % *KEYPHRASES *PHRASES
% % ----------------------
	% \label{LIFE-PHRASES}

	% \begin{compactdesc}
		% \item[] \textbf{\#times} | list
			% % \begin{compactdesc}
				% % \item[Bible anchor verses] \phantom{.}
					% % \begin{compactdesc}
						% % \item[Romans 5:5] \gr{}
					% % \end{compactdesc}
				% % \item[Commentary] ---
			% % \end{compactdesc}
	% \end{compactdesc}
	
% % ----------------------
% \subsection{Bible Dictionaries} % *DICTIONARIES
% % ----------------------
	% \label{LIFE-BD}

% % ----------------------
% \subsection{Restoration Usage} % *RESTORATION
% % ----------------------
	% \label{LIFE-RESTORATION}

	% % -----
	% \subsubsection{Joseph Smith} % *JS
	% % -----
		% \label{LIFE-JS}
	
		% \begin{compactdesc}
			% \item[{\hyperref[KEY]{Date/Source}}]
		% \end{compactdesc}
	
	% % -----
	% \subsubsection{Brigham Young} % *BY
	% % -----
		% \label{LIFE-BY}
	
		% \begin{compactdesc}
			% \item[{\hyperref[KEY]{Date/Source}}]
		% \end{compactdesc}
	
% % ----------------------
% \subsection{Commentary and Studies} % *COMMENTARY *STUDIES
% % ----------------------
	% \label{LIFE-COMMENTARY}

	% % -----
	% \subsubsection{Key Points}
	% % -----
		% \label{LIFE-KEYS}
	
		% \begin{compactitem}
			% \item
		% \end{compactitem}
		
	% % -----
	% \subsubsection{Sample Study}
	% % -----
		% \label{LIFE-study}